\section{The finite complex}
\label{sec:complex}

We specify the fixed triple system and certify the properties used in
both the kernel construction and the singular-tail bound.

Let $I=\{0,\ldots,12\}$ and $J=\{0,\ldots,21\}$. Associate to
each $j\in J$ the three-element subset of $I$ in Table~\ref{tab:complex};
we call these subsets \emph{faces} and also denote them by $j$.
Let $E$ be the set of unordered point pairs contained in a face.
The graph used throughout the paper is the incidence graph
\[
 H=(I\sqcup J,\mathscr C),\qquad
 \mathscr C=\{(i,j):i\in j\}.
\]
The copies of $I$ and $J$ in this disjoint union are distinct even when
their numerical labels agree. We call an incidence $(i,j)$ a
\emph{corner}. Thus $H$ has $35$ vertices and $66$ edges.

\begin{table}[htbp]
\centering
\small
\setlength{\tabcolsep}{7pt}
\begin{tabular}{c@{\qquad}ccc@{\qquad}ccc@{\qquad}c}
\toprule
Face&\multicolumn{3}{c}{Points}&\multicolumn{3}{c}{Opposite face}&Bits\\
 $j$&$i_1$&$i_2$&$i_3$&$j_{12}$&$j_{13}$&$j_{23}$&\\
\midrule
0 &0&1&3   &1&2&4   &000\\
1 &0&1&9   &0&3&7   &010\\
2 &0&2&3   &3&0&8   &001\\
3 &0&2&9   &2&1&9   &110\\
4 &1&3&10  &0&5&11  &001\\
5 &1&7&10  &6&4&16  &000\\
6 &1&7&12  &5&7&18  &100\\
7 &1&9&12  &1&6&21  &110\\
8 &2&3&4   &2&9&10  &101\\
9 &2&4&9   &8&3&12  &100\\
10&3&4&11  &8&11&13 &111\\
11&3&10&11 &4&10&17 &101\\
12&4&8&9   &13&9&21 &101\\
13&4&8&11  &12&10&19&011\\
14&5&6&7   &15&16&18&000\\
15&5&6&11  &14&17&19&110\\
16&5&7&10  &14&17&5 &010\\
17&5&10&11 &16&15&11&111\\
18&6&7&12  &14&20&6 &001\\
19&6&8&11  &20&15&13&010\\
20&6&8&12  &19&18&21&011\\
21&8&9&12  &12&20&7 &111\\
\bottomrule
\end{tabular}
\caption{The faces of the complex. The entry $j_{ab}$ is the other
face containing $\{i_a,i_b\}$. Bit positions are read from left to right.}
\label{tab:complex}
\end{table}

Join two faces when they share a pair. This defines the
\emph{face-neighbor graph}. For a pair $e\in E$, define
\begin{equation}
 a_e=\frac{1}{3}\#\{\text{bit positions at which the two faces on $e$
 differ}\}.
 \label{eq:complex-separation}
\end{equation}

Color refinement starts with a separate color for each part of $H$.
A refinement step replaces a vertex's color by the ordered pair
consisting of its old color and the histogram of its neighbors' old
colors. This is the classical color-refinement procedure; see
Weisfeiler and Leman~\cite{weisfeilerleman1968} and
McKay~\cite[Sections~2.3--2.6]{mckay1981} for its equitable-partition interpretation.

\begin{proposition}
\label{prop:complex}
The complex has the following properties.
\begin{enumerate}
\item Every pair in $E$ lies in exactly two faces, $|E|=33$, and the
face-neighbor graph is connected.
\item For each bit position, each of its two classes consists of
$11$ faces covering all $13$ points. Each class has an ordering in
which every face after the first shares a pair with an earlier face
and introduces exactly one new point.
\item Every $a_e$ is at least $1/3$, and some $a_e$ is greater than
$1/3$.
\item Color refinement of $H$, starting with a separate color for each
part, eventually gives distinct colors to all vertices within either
part.
\end{enumerate}
\end{proposition}

\begin{proof}
The neighbor entries in Table~\ref{tab:complex} pair the $66$
face--pair incidences into $33$ reciprocal pairs. Comparing the
triples in each indicated pair shows that their common point pair is
exactly the one specified. No other face contains that pair, as is
also checked by the same table.

Table~\ref{tab:complex-orders} gives all six orders and their
new-point certificates for the second assertion.

\begin{table}[htbp]
\centering
\small
\setlength{\tabcolsep}{5pt}
\begin{tabular}{cc@{\qquad}ll}
\toprule
\multicolumn{2}{c}{Bit class}&Face order&New points after the first face\\
Position&Value&&\\
\midrule
1&0&$0,1,2,4,5,16,14,18,20,19,13$
 &$9,2,10,7,5,6,12,8,11,4$\\
1&1&$3,9,8,10,11,17,15,12,21,7,6$
 &$4,3,11,10,5,6,8,12,1,7$\\
\addlinespace
2&0&$0,2,8,9,12,4,11,5,6,18,14$
 &$2,4,9,8,10,11,7,12,6,5$\\
2&1&$1,3,7,21,20,19,13,10,15,17,16$
 &$2,12,8,6,11,4,3,5,10,7$\\
\addlinespace
3&0&$0,1,3,9,7,6,5,16,14,15,19$
 &$9,2,4,12,7,10,5,6,11,8$\\
3&1&$2,8,10,11,4,17,13,12,21,20,18$
 &$4,11,10,1,5,8,9,12,6,7$\\
\bottomrule
\end{tabular}
\caption{An exposure order for each bit class. Every face after the
first shares a pair with an earlier face and introduces the
corresponding point in the last column. Bit positions are read from
left to right in Table~\ref{tab:complex}.}
\label{tab:complex-orders}
\end{table}

The root face contributes its three points, and these ten new points
give all of $I$ in each case. The shared-pair condition follows from
the neighbor columns of Table~\ref{tab:complex}. Thus every bit class
is connected. The first two classes are joined, for example, by
faces $2$ and $3$, so the whole face-neighbor graph is connected.
Inspection of the bit labels gives separation in at least one
position on every pair. The pair $\{0,2\}$ has adjacent labels
$001$ and $110$, so in fact $a_{\{0,2\}}=1$.

It remains to verify the color-refinement assertion. First consider
the point graph $(I,E)$. Its point degrees $d_i$ and the numbers
$n_4(i),n_5(i),n_6(i)$ of neighbors of degrees $4,5,6$ are as follows:
\begin{equation}
\begin{array}{c|c|ccc@{\qquad}c|c|ccc}
i&d_i&n_4&n_5&n_6&i&d_i&n_4&n_5&n_6\\\hline
0&4&1&0&3&7&5&1&3&1\\
1&6&1&3&2&8&5&0&3&2\\
2&4&1&1&2&9&6&2&3&1\\
3&6&2&2&2&10&5&1&1&3\\
4&5&1&1&3&11&6&1&4&1\\
5&4&0&3&1&12&5&0&3&2\\
6&5&1&3&1&&&&&
\end{array}
\label{eq:complex-colors}
\end{equation}
Only the pairs $\{4,10\}$, $\{6,7\}$, and $\{8,12\}$ have
equal signatures in Equation~\eqref{eq:complex-colors}. The first is
separated by adjacency to the already distinguished point $2$:
$4$ is adjacent to $2$, whereas $10$ is not. The other two are
separated by adjacency to the already distinguished point $11$:
$6$ and $8$ are adjacent to $11$, whereas $7$ and $12$ are not.
Hence refinement on the point graph is discrete.

This point refinement can be recovered from incidence refinement.
The degree of $i$ in $H$ equals $d_i$: counting other points in its
incident faces gives twice each neighbor in $(I,E)$. More generally,
suppose the current incidence color of every point determines a point
coloring \(\ell\) through a fixed projection. One further incidence
step determines, at each face \(j\), the color-indexed histogram
\[
 m_j(\gamma)=\#\{k\in j:\ell(k)=\gamma\}.
\]
The next step determines the multiset of these histograms over the
faces incident to each point. For every color
\(\gamma\), this gives
\[
 \sum_{j\ni i}
 \left(m_j(\gamma)-\1_{\{\gamma=\ell(i)\}}\right)
 =2\,\#\{k:\{i,k\}\in E,\ \ell(k)=\gamma\}.
\]
The subtraction removes only the occurrence of \(i\) in each face,
even when other points have the same color. Each neighbor is counted
twice because its pair with \(i\) lies in two faces. Dividing by two
therefore recovers the point neighbor-color histogram.
Thus two incidence steps recover each further point refinement step.
After the point colors are discrete, one more step distinguishes all
faces, since their three-element point sets are distinct. This proves
the last assertion without any assumption about a chosen refinement
implementation.
\end{proof}
